% Part: first-order-logic
% Chapter: natural-deduction
% Section: soundness-identity

\documentclass[../../../include/open-logic-section]{subfiles}

\begin{document}

\olfileid{fol}{ntd}{sid}

\olsection{\usetoken{S}{identity}తో నిర్దుష్టత}

\begin{prop}
$\eq$కు నియమాలు గల సహజ నిగమనం నిర్దుష్టమైనది.
\end{prop}

\begin{proof}
$\eq[t][t]$ రూపంలోని ప్రతి !!{formula} చెల్లుబాటు అవుతుంది; ఎందుకంటే
ప్రతి !!{structure}~$\Struct M$కు $\Sat{M}{\eq[t][t]}$. ($t$ సంవృత
పదం, అంటే అది ఏ చరరాశులనూ కలిగి ఉండదని అనుకుంటున్నాం; కనుక చరరాశి
నిర్దేశాలు అప్రస్తుతం.)

ఒక !!a{derivation}లో చివరి నిగమనం \Elim{\eq} అనుకుందాం; అంటే
!!{derivation} కింది రూపంలో ఉంటుంది:
\begin{prooftree}
  \AxiomC{$\Gamma_1$}
  \RightLabel{$\delta_1$}
  \DeduceC{$\eq[t_1][t_2]$}
  \AxiomC{$\Gamma_2$}
  \RightLabel{$\delta_2$}
  \DeduceC{$!A(t_1)$}
  \RightLabel{\Elim{\eq}}
  \BinaryInfC{$!A(t_2)$}
\end{prooftree}
పూర్వాధారాలు $\eq[t_1][t_2]$, $!A(t_1)$లు వరుసగా !!{undischarged}
పరికల్పనలు~$\Gamma_1$, $\Gamma_2$ల నుంచి !!{derive}d అయ్యాయి.
$\Gamma_1 \cup \Gamma_2$ నుంచి $!A(t_2)$ అనుగమిస్తుందని చూపాలి.
$\Sat{M}{\Gamma_1 \cup \Gamma_2}$ అయ్యే ఒక
!!a{structure}~$\Struct{M}$ను పరిశీలించండి. ఆగమన పరికల్పన ప్రకారం
$\Sat{M}{!A(t_1)}$, $\Sat{M}{\eq[t_1][t_2]}$. అందువల్ల
$\Value{t_1}{M} = \Value{t_2}{M}$. $s$ ఏదైనా చరరాశి నిర్దేశం అనీ,
$m = \Value{t_1}{M} = \Value{t_2}{M}$ అనీ అందాం.
\olref[fol][syn][ext]{prop:ext-formulas} ప్రకారం
$\Sat{M}{!A(t_1)}[s]$ అయితే, అప్పుడు మాత్రమే
$\Sat{M}{!A(x)}[\Subst{s}{m}{x}]$; అలాగే అప్పుడు మాత్రమే
$\Sat{M}{!A(t_2)}[s]$. $\Sat{M}{!A(t_1)}$ కనుక
$\Sat{M}{!A(t_2)}$.
\end{proof}

\end{document}
